\documentclass{article}
\usepackage{apstats-poster}
% Formula IDs: zscore, empirical-rule
\colorlet{PosterAccent}{UnitBlue}

\begin{document}
\begin{posterpage}
\PosterHeader[86pt]{UnitBlue}{2}{PROBABILITY, RANDOM VARIABLES \& SAMPLING DISTRIBUTIONS}
  {P07}{THE NORMAL DISTRIBUTION}
  {new 2.11 \quad\textbullet\quad old 1.10}

\PosterZone{4.72in}{16.40in}{ZONE A · THE CURVE + THE FORMULA}{%
  \begin{center}
    \begin{tikzpicture}[x=2.42in,y=7in]
      \fill[UnitBlue!15] (-3,0) -- plot[domain=-3:3,samples=100]
        (\x,{.39894*exp(-\x*\x/2)}) -- (3,0) -- cycle;
      \fill[UnitBlue!35] (-2,0) -- plot[domain=-2:2,samples=80]
        (\x,{.39894*exp(-\x*\x/2)}) -- (2,0) -- cycle;
      \fill[UnitBlue!65] (-1,0) -- plot[domain=-1:1,samples=60]
        (\x,{.39894*exp(-\x*\x/2)}) -- (1,0) -- cycle;
      \draw[PosterInk,line width=3pt,-{Stealth[length=12pt]}] (-3.6,0)--(3.6,0);
      \draw[UnitBlue,line width=4pt] plot[domain=-3.5:3.5,samples=120]
        (\x,{.39894*exp(-\x*\x/2)});
      \foreach \v in {-3,-2,-1,0,1,2,3} {
        \draw[PosterInk,line width=2pt] (\v,-.008)--(\v,.008);
        \draw[white,line width=1.5pt,dashed] (\v,0)--(\v,{.39894*exp(-\v*\v/2)});
      }
      \foreach \v/\label in {-3/{\mu-3\sigma},-2/{\mu-2\sigma},-1/{\mu-\sigma},0/{\mu},1/{\mu+\sigma},2/{\mu+2\sigma},3/{\mu+3\sigma}}
        \node[anchor=north,font=\fontsize{25pt}{29pt}\selectfont] at (\v,-.018) {$\label$};
      \draw[UnitBlue,line width=3pt,{Stealth[length=9pt]}-{Stealth[length=9pt]}] (-1,-.13)--(1,-.13);
      \node[anchor=center,fill=white,font=\bfseries\fontsize{29pt}{33pt}\selectfont] at (0,-.13) {68\% within \(1\sigma\)};
      \draw[UnitBlue!70,line width=3pt,{Stealth[length=9pt]}-{Stealth[length=9pt]}] (-2,-.23)--(2,-.23);
      \node[anchor=center,fill=white,font=\bfseries\fontsize{29pt}{33pt}\selectfont] at (0,-.23) {95\% within \(2\sigma\)};
      \draw[UnitBlue!45,line width=3pt,{Stealth[length=9pt]}-{Stealth[length=9pt]}] (-3,-.33)--(3,-.33);
      \node[anchor=center,fill=white,font=\bfseries\fontsize{29pt}{33pt}\selectfont] at (0,-.33) {99.7\% within \(3\sigma\)};
    \end{tikzpicture}
  \end{center}
  \vspace{.08in}
  \begin{minipage}[c]{8.0in}
    \FormulaCardSized{71pt}{z=\frac{x-\mu}{\sigma}}
  \end{minipage}\hfill
  \begin{minipage}[c]{11.5in}
    \TextCard{\centering\fontsize{37pt}{50pt}\selectfont
      \(x\ \longrightarrow\ z\ \longrightarrow\ \text{area}\)\par
      {\fontsize{25pt}{31pt}\selectfont normalcdf: value to probability}\par\vspace{.16in}
      \(\text{area}\ \longrightarrow\ z\ \longrightarrow\ x\)\par
      {\fontsize{25pt}{31pt}\selectfont invNorm: probability to value}}
  \end{minipage}
}

\PosterZone{16.55in}{19.85in}{ZONE B · WHAT THE SYMBOLS MEAN}{%
  \begin{minipage}[t]{9.8in}
    \MiniHeading{CENTER + SPREAD}\par\vspace{.10in}
    {\fontsize{32pt}{41pt}\selectfont
      \(X\sim N(\mu,\sigma)\)\par
      \(\mu\): mean \qquad \(\sigma\): standard deviation\par
      Total area under the curve is 1.}
  \end{minipage}\hfill
  \begin{minipage}[t]{9.8in}
    \MiniHeading{STANDARDIZE THE DISTANCE}\par\vspace{.10in}
    {\fontsize{32pt}{41pt}\selectfont
      \(z\) counts SDs from the mean.\par
      Positive: above \(\mu\). Negative: below \(\mu\).\par
      At the mean, \(z=0\).}
  \end{minipage}
}

\PosterZone{20.00in}{23.75in}{ZONE C · SAY IT IN WORDS}{%
  \SentenceFrame{“\PosterBlank{1.5in} is \PosterBlank{1.0in} standard deviations
    [above / below] the mean of \PosterBlank{3.3in}.”}
  \vspace{.18in}
  \SentenceFrame{“About \PosterBlank{1.0in}\% of \PosterBlank{2.4in} have
    \PosterBlank{2.8in} below \PosterBlank{1.3in}.”}
}

\PosterZone{23.90in}{27.45in}{ZONE D · WATCH OUT}{%
  \begin{minipage}[t]{9.7in}
    \MiniHeading{DRAW BEFORE YOU CALCULATE}\par\vspace{.16in}
    {\fontsize{33pt}{42pt}\selectfont
      Draw the curve. Mark the value.\par
      Shade the requested region first.}
  \end{minipage}\hfill
  \begin{minipage}[t]{9.7in}
    \MiniHeading{PERCENTILE = AREA TO THE LEFT}\par\vspace{.16in}
    {\fontsize{31pt}{40pt}\selectfont
      Multiply the left-tail area by 100.\par
      The 68--95--99.7 rule gives approximate percentages for a normal model.}
  \end{minipage}
}
\WorksheetTag{u1\_l10}
\end{posterpage}
\end{document}
